
\documentclass[a4paper, 11pt]{article}

% ── Fonts ──
\usepackage{fontspec}
\setmainfont{Arial}
\setsansfont{Arial}
\setmonofont{Consolas}

% ── Math ──
\usepackage{amsmath, amssymb, amsfonts}

% ── Math macros (matching the book) ──
\newcommand{\R}{\mathbb{R}}
\newcommand{\N}{\mathbb{N}}
\newcommand{\Z}{\mathbb{Z}}
\newcommand{\Q}{\mathbb{Q}}
\newcommand{\C}{\mathbb{C}}
\newcommand{\F}{\mathbb{F}}
\newcommand{\M}{\mathbb{M}}
\newcommand{\Set}[1]{\left\{#1\right\}}

% ── Page layout ──
\usepackage[margin=2cm]{geometry}

% ── Colors for environment boxes ──
\usepackage{tcolorbox}
\tcbuselibrary{breakable, skins}

\definecolor{defcolor}{RGB}{220, 53, 69}
\definecolor{thmcolor}{RGB}{230, 126, 34}

% ── RTL support (must be loaded LAST) ──
\usepackage{bidi}

\newtcolorbox{envbox}[2][]{%
  enhanced, breakable,
  colback=#1!3,
  colframe=#1!60,
  coltitle=black,
  fonttitle=\bfseries,
  title={#2},
  boxrule=0.6pt,
  arc=2pt,
  left=8pt, right=8pt, top=6pt, bottom=6pt,
  attach boxed title to top right={yshift=-2mm, xshift=-4mm},
  boxed title style={
    arc=2pt, boxrule=0.4pt,
    colback=#1!20,
    colframe=#1!60,
  },
  before skip=10pt,
  after skip=10pt,
}

% ── Section style ──
\usepackage{titlesec}
\titleformat{\section}{\Large\bfseries}{}{0pt}{}
\titlespacing*{\section}{0pt}{20pt}{10pt}
\setcounter{secnumdepth}{0}

% ── Hebrew TOC title ──
\renewcommand{\contentsname}{תוכן עניינים}

% ── No paragraph indent, add spacing ──
\setlength{\parindent}{0pt}
\setlength{\parskip}{6pt}

\begin{document}

% ── Title Page ──
\vspace*{2cm}
\begin{center}
{\Huge\bfseries סיכום הגדרות ומשפטים}\\[14pt]
{\Large אלגברה לינארית להנדסה}\\[10pt]
{\normalsize ד״ר יונתן שלח \quad·\quad ד״ר נדב מאיר}
\end{center}
\vspace{1.5cm}

\tableofcontents
\newpage

\section{פרק 0: קבוצות ומספרים}

\begin{envbox}[defcolor]{הגדרה 0.2}
קבוצה היא אוסף כלשהו של איברים (לא בהכרח מספרים) ללא חשיבות לסדר הופעתם.
\end{envbox}

\begin{envbox}[defcolor]{הגדרה 0.5}
עבור $z=x+yi$ מרוכב עם $x,y\in\R$ נגדיר את החלק הממשי $\text{Re}(z)=x$ ואת החלק המדומה $\text{Im}(z)=y$. בנוסף, נגדיר את המספר הצמוד $\overline{z}=x-yi$.
\end{envbox}

\begin{envbox}[defcolor]{הגדרה 0.8}
נאמר ש-$A$ מוכלת ב-$B$ ונסמן $A\subseteq B$ אם כל איבר של $A$ הוא גם איבר של $B$.

באופן קצת יותר מתמטי: $A\subseteq B$ אם לכל $a\in A$ מתקיים $a\in B$. 

אם ההיפך הוא הנכון, כלומר קיים $a\in A$ עבורו $a\notin B$, אז נסמן $A\nsubseteq B$ ונאמר ש-$A$ לא מוכלת ב-$B$.
\end{envbox}

\begin{envbox}[defcolor]{הגדרה 0.11}
נאמר ששתי קבוצות $A,B$ הן שוות ונסמן $A=B$ אם יש בהן בדיוק אותם האיברים, כלומר מתקיים $x\in A$ אם ורק אם $x\in B$.
\end{envbox}

\begin{envbox}[defcolor]{הגדרה 0.12}
כדי להוכיח כי $A=B$, הדרך המקובלת היא להוכיח כי $A\subseteq B$ וגם $B\subseteq A$. דרך הוכחה זו נקראת הכלה דו-צדדית.
\end{envbox}

\begin{envbox}[thmcolor]{טענה 0.14}
$\N\subseteq\Z\subseteq\Q\subseteq\R\subseteq\C$ וכל הקבוצות שונות זו מזו.
\end{envbox}

\begin{envbox}[defcolor]{הגדרה 0.15}
בהינתן שתי קבוצות $A,B$ נגדיר את הקבוצות הבאות:

a. החיתוך $A\cap B$ הוא קבוצת כל האיברים השייכים גם ל-$A$ וגם ל-$B$.
b. האיחוד $A\cup B$ הוא קבוצת כל האיברים השייכים לפחות לאחת משתי הקבוצות $A,B$.
\end{envbox}

\begin{envbox}[thmcolor]{טענה 0.19}
לכל שתי קבוצות $A,B$ מתקיים

\[
A\cap B\subseteq A\subseteq A\cup B.
\]
\end{envbox}

\newpage

\section{פרק 1: תכונות ושימושים של המספרים המרוכבים}

\begin{envbox}[thmcolor]{טענה 1.4}
לכל $z_{1},z_{2},z_{3}\in\mathbb{C}$ מתקיים:

a. חוקי החילוף לחיבור וכפל: $z_{1}+z_{2}=z_{2}+z_{1}$ וגם $z_{1}z_{2}=z_{2}z_{1}$.
b. חוקי הקיבוץ לחיבור וכפל: $(z_{1}+z_{2})+z_{3}=z_{1}+(z_{2}+z_{3})$ וגם $(z_{1}z_{2})z_{3}=z_{1}(z_{2}z_{3})$.
c. חוק הפילוג: $(z_{1}+z_{2})z_{3}=z_{1}z_{3}+z_{2}z_{3}$
d. $0$ נייטרלי ביחס לחיבור: $z+0=z$ לכל $z\in\mathbb{C}$.
e. $1$ נייטרלי ביחס לכפל: $z\cdot1=z$ לכל $z\in\mathbb{C}$.
\end{envbox}

\begin{envbox}[thmcolor]{טענה 1.6}
לכל $z_{1},z_{2}\in\mathbb{C}$ מתקיים:

a. $\overline{z_{1}+z_{2}}=\overline{z_{1}}+\overline{z_{2}}$
b. $\overline{z_{1}z_{2}}=\overline{z_{1}}\cdot\overline{z_{2}}$
\end{envbox}

\begin{envbox}[thmcolor]{טענה 1.7}
לכל $z\in\mathbb{C}$ מתקיים:

a. $z+\overline{z}=2\text{Re}(z)$
b. $z-\overline{z}=2i\text{Im}(z)$
c. $z\overline{z}\in\mathbb{R}$
\end{envbox}

\begin{envbox}[defcolor]{הגדרה 1.12}
נגדיר

\[
e^{i\theta}=\cos\theta+i\sin\theta.
\]
\end{envbox}

\begin{envbox}[thmcolor]{משפט 1.13}
לכל $\theta\in\mathbb{R},\,n\in\mathbb{Z}$ מתקיים

\[
(\cos\theta+i\sin\theta)^{n}=\cos(n\theta)+i\sin(n\theta)
\]
.
\end{envbox}

\begin{envbox}[thmcolor]{טענה 1.14}
לכל $\alpha,\beta\in\mathbb{R}$ מתקיים

\[
e^{i\alpha}e^{i\beta}=e^{i(\alpha+\beta)}
\]
.
\end{envbox}

\begin{envbox}[thmcolor]{טענה 1.18}
יהי $w\in\mathbb{C}$ מספר נתון השונה מ-$0$, ונניח כי $w=re^{i\theta}$ בהצגה פולרית. אז למשוואה $z^{n}=w$ יש בדיוק $n$ שורשים (פתרונות) מרוכבים בנעלם $z$ הנתונים ע״י

\[
z_{k}=\sqrt[n]{r}e^{i\frac{\theta+2\pi k}{n}}.
\]

עבור $k\in\Set{0,1,...,n-1}$.
\end{envbox}

\newpage

\section{פרק 2: גיאומטריה במישור ובמרחב}

\begin{envbox}[defcolor]{הגדרה 2.1}
המישור $\mathbb{R}^2$ הוא קבוצת כל הנקודות בעלות שתי קוארדינטות ממשיות, כלומר 
\[
.\mathbb{R}^2=\Set{(x,y)|x,y\in\mathbb{R}}.
\]
\end{envbox}

\begin{envbox}[defcolor]{הגדרה 2.3}
וקטור $\vec{v}=(x,y)$ במישור הוא זוג סדור של מספרים ממשיים, כלומר איבר של $\mathbb{R}^2$.
\end{envbox}

\begin{envbox}[thmcolor]{טענה 2.12}
$\|c\cdot(a,b)\| = |c|\cdot\|(a,b)\|$ לכל וקטור $(a,b)\in\mathbb{R}^2$ ולכל סקלר $c\in\mathbb{R}$.
\end{envbox}

\begin{envbox}[defcolor]{הגדרה 2.18}
עבור וקטורים $(a_1,a_2),(b_1,b_2)\in\mathbb{R}^2$ נגדיר  
$(a_1,a_2)\cdot(b_1,b_2)=a_1b_1+a_2b_2$.
\end{envbox}

\begin{envbox}[thmcolor]{טענה 2.20}
יהיו $\vec{v_1}=(x_1,y_1),\,\vec{v_2}=(x_2,y_2),\,\vec{v_3}=(x_3,y_3)\in\mathbb{R}^2$ וקטורים. אז מתקיים:

1. $\vec{v_1}\cdot\vec{v_2}=\vec{v_2}\cdot\vec{v_1}$  
2. $(\vec{v_1}+\vec{v_3})\cdot\vec{v_2}=\vec{v_1}\cdot\vec{v_2}+\vec{v_3}\cdot\vec{v_2}$  
3. $\vec{v_1}\cdot\vec{v_1}=\|\vec{v_1}\|^2$  
4. $\vec{v_1}\cdot(c\vec{v_2})=c(\vec{v_1}\cdot\vec{v_2})$ לכל סקלר $c\in\mathbb{R}$  
5. $\|\vec{v_1}\|=0\iff\vec{v_1}=(0,0)$  
6. $\|\frac{\vec{v_1}}{\|\vec{v_1}\|}\|=1$ אם $\vec{v_1}\neq(0,0)$
\end{envbox}

\begin{envbox}[thmcolor]{טענה 2.21}
יהיו $\vec{v},\vec{w}$ שני וקטורים ב-$\mathbb{R}^2$ השונים מ-$(0,0)$, ותהי $\alpha$ הזווית הקטנה ביניהם. אז מתקיים:

\[
\cos\alpha = \frac{\vec{v}\cdot\vec{w}}{\|\vec{v}\|\cdot\|\vec{w}\|}.
\]
 {\#eq-cosine-scalar}
\end{envbox}

\begin{envbox}[thmcolor]{מסקנה 2.22}
וקטורים $\vec{v},\vec{w}$ השונים מ-$(0,0)$ הם מאונכים/ניצבים אם ורק אם $\vec{v}\cdot\vec{w}=0$. במקרה זה נסמן $\vec{v}\perp\vec{w}$.
\end{envbox}

\begin{envbox}[defcolor]{הגדרה 2.26}
וקטור המאונך לכל וקטור שמתאים לשתי נקודות על ישר (התחלה וסיום) נקרא נורמל לישר.
\end{envbox}

\begin{envbox}[defcolor]{הגדרה 2.32}
שני וקטורים $\vec{v},\vec{w}\in\mathbb{R}^2$ נקראים קו-לינאריים אם קיים $c\in\mathbb{R}$ כך ש-$\vec{v}=c\vec{w}$ או $\vec{w}=c\vec{v}$.
\end{envbox}

\begin{envbox}[defcolor]{הגדרה 2.42}
המרחב $\mathbb{R}^3$ הוא קבוצת כל הנקודות בעלות שלוש קוארדינטות ממשיות, כלומר 
\[
.\mathbb{R}^3=\Set{(x,y,z)|x,y,z\in\mathbb{R}}.
\]
\end{envbox}

\begin{envbox}[defcolor]{הגדרה 2.45}
וקטור $\vec{v}=(x,y,z)$ במרחב הוא שלשה סדורה של מספרים ממשיים, כלומר איבר ב- $.\mathbb{R}^3$
\end{envbox}

\begin{envbox}[thmcolor]{טענה 2.53}
$\|c\cdot(x,y,z)\|=|c|\cdot\|(x,y,z)\|.$ לכל וקטור $(x,y,z)\in\mathbb{R}^3$ ולכל סקלר $c\in\mathbb{R}$.
\end{envbox}

\begin{envbox}[defcolor]{הגדרה 2.55}
עבור וקטורים $(a_1,a_2,a_3),(b_1,b_2.b_3)\in\mathbb{R}^3$ נגדיר 
\[
.(a_1,a_2,a_3)\cdot(b_1,b_2,b_3)=a_1b_1+a_2b_2+a_3b_3.
\]
\end{envbox}

\begin{envbox}[thmcolor]{טענה 2.57}
יהיו $\vec{v_1}=(x_1,y_1\text{,}z_{!}),\,\vec{v_2}=(x_2,y_2,z_2),\,\vec{v_3}=(x_3,y_3,z_3)\in\mathbb{R}^3$ וקטורים. אז מתקיים:

א. $\vec{v_1}\cdot\vec{v_2}=\vec{v_2}\cdot\vec{v_1}$.

ב. $(\vec{v_1}+\vec{v_3})\cdot\vec{v_2}=\vec{v_2}\cdot\vec{v_1}$.

ג. $\vec{v_1}\cdot\vec{v_1}=\|\vec{v_1}\|^2$.

ד. $\vec{v_1}\cdot(c\vec{v_2})=c(\vec{v_1}\cdot\vec{v_2})$.

ה. ${\|\vec{v_1}\|}=0\iff\vec{v_1}=(0,0)$.

ו. $\|\frac{\vec{v_1}}{\|\vec{v_1}\|}\|=1$ אם $\vec{v_1}\neq(0,0)$.
\end{envbox}

\begin{envbox}[thmcolor]{טענה 2.58}
יהיו $\vec{v},\vec{w}$ שני וקטורים ב-$\mathbb{R}^3$ השונים מ-$(0,0,0)$, ותהי $\alpha$ הזווית הקטנה ביניהם. אז מתקיים 
\[
.\cos\alpha=\frac{\vec{v}\cdot\vec{w}}{\|\vec{v}\|\cdot\|\vec{w}\|}.
\]
\end{envbox}

\begin{envbox}[thmcolor]{מסקנה 2.60}
נניח כי $\vec{v},\vec{w}$ שני וקטורים ב-$\mathbb{R}^3$ השונים מ-$(0,0,0)$. אז הם מאונכים זה לזה אם ורק אם $\vec{v}\cdot\vec{w}=0$.
\end{envbox}

\newpage

\section{פרק 3: מערכות משוואות לינאריות}

\begin{envbox}[defcolor]{הגדרה 3.1}
משוואה לינארית במשתנים $x_1,x_{2,}...,x_n$ היא משוואה מהצורה 

\[
a_1x_1 + a_2x_2 + a_3x_3 + \cdots + a_n x_n = b
\]
 {\#eq-linear-eq}

כאשר $a_1,...,a_n$ הם סקלרים נתונים שנקראים מקדמים, ו-$b$ הוא סקלר נתון שנקרא קבוע או איבר חופשי. אם כל הסקלרים והמשתנים ממשיים, נאמר שזו משוואה לינארית מעל $\mathbb R$. אם לפחות סקלר אחד הוא מרוכב ו/או המשתנים מרוכבים, נאמר שזו משוואה לינארית מעל $\mathbb C$. אם נרצה להכליל בלי לציין שמדובר ב-$\mathbb R$ או $\mathbb C$, נכתוב $\mathbb F$ ונזכור שקבוצה זו היא או $\mathbb R$ או $\mathbb C$.
\end{envbox}

\begin{envbox}[defcolor]{הגדרה 3.3}
סדרת סקלרים $(c_1,c_2,...,c_n)$ נקראת $n$-יה, או וקטור באורך $n$. נגדיר את מרחב ה-$n$-יות מעל $\mathbb R$: 
\[
\mathbb R^{n}=\Set{(c_1,c_2,....,c_n)|c_1,c_2,....,c_n\in\mathbb R}.
\]
 וגם את מרחב ה-$n$-יות מעל $\mathbb C$: 
\[
\mathbb C^{n}=\Set{(c_1,c_2,....,c_n)|c_1,c_2,....,c_n\in\mathbb C}.
\]

לעיתים נכתוב $\mathbb F^n$ במקום לפרט אם הכוונה היא ל-$\mathbb R^n$ או ל-$\mathbb C^n$.
\end{envbox}

\begin{envbox}[defcolor]{הגדרה 3.4}
פתרון של משוואה לינארית הינו $n$-יה $(c_1,c_2,...,c_n)\in\mathbb F^{n}$ כך שבהצבת 
\[
x_1=c_1,x_2=c_2,...,x_n=c_n.
\]

המשוואה מתקיימת.
\end{envbox}

\begin{envbox}[defcolor]{הגדרה 3.9}
מערכת משוואות לינאריות (ממ״ל) במשתנים $x_1,x_{2,}...,x_n$ היא אוסף של מספר כלשהו ($m$) של משוואות לינאריות מהצורה

\[
\begin{cases}
\begin{alignedat}{7}
a_{11}x_1 &{}+{}& a_{12}x_2 &{}+{}& a_{13}x_3 &{}+{}& \dots + a_{1n}x_n &= b_1 \\[4pt]
a_{21}x_1 &{}+{}& a_{22}x_2 &{}+{}& a_{23}x_3 &{}+{}& \dots + a_{2n}x_n &= b_2 \\[4pt]
a_{31}x_1 &{}+{}& a_{32}x_2 &{}+{}& a_{33}x_3 &{}+{}& \dots + a_{3n}x_n &= b_3 \\[4pt]
\vdots& & & & \vdots & & & \;\;\vdots \\[4pt]
a_{m1}x_1 &{}+{}& a_{m2}x_2 &{}+{}& a_{m3}x_3 &{}+{}& \dots + a_{mn}x_n &= b_m
\end{alignedat}
\end{cases}
\]
 {\#eq-linear-eqs}

אם כל המקדמים והקבועים הם מספרים ממשיים, נאמר שהממ״ל מעל $\mathbb R$. אחרת, נאמר שהממ״ל מעל $\mathbb C$.
\end{envbox}

\begin{envbox}[defcolor]{הגדרה 3.10}
פתרון של ממ״ל הינו סדרת סקלרים $(c_1,c_2,...,c_n)$ כך שבהצבת 
\[
x_1=c_1,x_2=c_2,...,x_n=c_n
\]
 כל המשוואות בממ״ל מתקיימות.
\end{envbox}

\begin{envbox}[defcolor]{הגדרה 3.14}
מטריצה מסדר $m\times n$ מעל $\mathbb F$ היא טבלה בת $m$ שורות ו-$n$ עמודות כך שבכל משבצת מופיע סקלר ב-$\mathbb F$. בהקשר של הממ״ל 
\[
,\begin{cases}
\begin{alignedat}{7}
a_{11}x_1 &{}+{}& a_{12}x_2 &{}+{}& a_{13}x_3 &{}+{}& \dots + a_{1n}x_n &= b_1 \\[4pt]
a_{21}x_1 &{}+{}& a_{22}x_2 &{}+{}& a_{23}x_3 &{}+{}& \dots + a_{2n}x_n &= b_2 \\[4pt]
a_{31}x_1 &{}+{}& a_{32}x_2 &{}+{}& a_{33}x_3 &{}+{}& \dots + a_{3n}x_n &= b_3 \\[4pt]
\vdots& & & & \vdots & & & \;\;\vdots \\[4pt]
a_{m1}x_1 &{}+{}& a_{m2}x_2 &{}+{}& a_{m3}x_3 &{}+{}& \dots + a_{mn}x_n &= b_m
\end{alignedat}
\end{cases}
\]
 נגדיר את מטריצת המקדמים (שנקראת גם מטריצת מקדמים מצומצמת):

\[
A=\begin{pmatrix}a_{11} & a_{12} & \cdots & a_{1n}\\
a_{21} & a_{22} & \cdots & a_{2n}\\
\vdots & \vdots & \ddots & \vdots\\
a_{m1} & a_{m2} & \cdots & a_{mn}
\end{pmatrix}
\]

נגדיר גם את מטריצת המקדמים המורחבת מסדר $m\times(n+1)$ שכוללת עמודה בצד ימין (אחרי קו מפריד) בשביל הקבועים:

\[
\left( \begin{array}{cccc|c}
    a_{11} & a_{12} & \cdots & a_{1n} & b_1 \\
    a_{21} & a_{22} & \cdots & a_{2n} & b_2 \\
    \vdots & \vdots & \ddots & \vdots & \vdots \\
    a_{m1} & a_{m2} & \cdots & a_{mn} & b_m
\end{array} \right)
\]

נגדיר בנוסף את וקטור המשתנים $x=\begin{pmatrix}x_1\\
x_2\\
\vdots\\
x_n
\end{pmatrix}$ וגם את וקטור הקבועים $b=\begin{pmatrix}b_1\\
b_2\\
\vdots\\
b_m
\end{pmatrix}$. 
דרך מקוצרת לכתיבת הממ״ל  היא המשוואה הוקטורית $Ax=b$.
\end{envbox}

\begin{envbox}[defcolor]{הגדרה 3.18}
מטריצה נקראת מדורגת אם:  

א. אם במטריצה יש שורות אפסים, אז הן מופיעות מתחת לכל שאר השורות.  

ב. בכל שורה שאינה שורת אפסים, האיבר הראשון (משמאל) שאינו $0$ נקרא האיבר המוביל של שורה זו, והוא מופיע משמאל לאיבר המוביל של השורה מתחתיו (אם קיימת כזו). בפרט, מתחת לכל איבר מוביל יש אפסים בלבד.
\end{envbox}

\begin{envbox}[defcolor]{הגדרה 3.20}
מטריצה נקראת מדורגת קנונית אם:

א. היא מדורגת, כלומר מקיימת את שני תנאי ההגדרה הקודמת.  

ב. האיבר המוביל בכל שורה (אם הוא קיים) הוא $1$. אפשר גם לקרוא לו אחד מוביל.  

ג. בכל עמודה של אחד מוביל, כל שאר האיברים הם $0$ (מתחת וגם מעל לאחד המוביל).
\end{envbox}

\begin{envbox}[defcolor]{הגדרה 3.22}
נניח שקיבלנו צורה מדורגת קנונית של מטריצת מקדמים מצומצמת $A$ של ממ״ל $Ax=b$.  

א. כל משתנה שמתאים לעמודה בה מופיע איבר מוביל נקרא משתנה תלוי. שאר המשתנים נקראים חופשיים.    

ב. למספר המשתנים התלויים נקרא הדרגה של הממ״ל/מטריצה. נסמן $\mathop{\mathrm{rank}}(A)$.
\end{envbox}

\begin{envbox}[defcolor]{הגדרה 3.25}
תהי

\[
A=\begin{pmatrix}a_{11} & a_{12} & \cdots & a_{1n}\\
a_{21} & a_{22} & \cdots & a_{2n}\\
\vdots & \vdots & \ddots & \vdots\\
a_{m1} & a_{m2} & \cdots & a_{mn}
\end{pmatrix}
\]
 מטריצה מסדר $m\times n$ מעל $\mathbb F$. נסמן ב-$R_1,R_2,...,R_m$ את שורות המטריצה $A$, כלומר $R_{i}=(a_{i1},a_{i2},...,a_{in})$ לכל $1\leq i\leq m$. להלן שלושה סוגי פעולות על שורות $A$ הנקראות פעולות שורה אלמנטריות:

א. כפל השורה $R_{i}$ בסקלר $c\in F$ שאינו $0$. נסמן פעולה זו כך: $cR_{i}\rightarrow R_{i}$.  

ב. הוספה לשורה $R_{i}$ את הכפולה של השורה $R_{j}$ בסקלר $c\in \mathbb F$ כאשר $i\neq j$. נסמן $R_{i}+cR_{j}\rightarrow R_{i}$. השורה $R_{j}$ נשארת ללא שינוי.  

ג. החלפת השורות $R_{i}$ ו-$R_{j}$ עבור $i\neq j$. נסמן $R_{i}\leftrightarrow R_{j}$.
\end{envbox}

\begin{envbox}[defcolor]{הגדרה 3.28}
מטריצות $A,A'$ נקראות שקולות שורה אם ניתן להגיע מ- $A$ ל- $A'$ ע״י סדרה סופית של פעולות שורה אלמנטריות.
\end{envbox}

\begin{envbox}[thmcolor]{טענה 3.30}
. נסתכל על הממ״ל $Ax=b$ עם מטריצת מקדמים מורחבת $(A|b)$, וגם על הממ״ל $A'x=b'$ עם מטריצת מקדמים מורחבת $(A'|b')$. אם $(A|b),(A'|b')$ הן מטריצות שקולות שורה, אז קבוצת הפתרונות של הממ״ל $Ax=b$ שווה לקבוצת הפתרונות של הממ״ל $A'x=b'$.
\end{envbox}

\begin{envbox}[defcolor]{הגדרה 3.32}
כאשר וקטור הקבועים של הממ״ל $Ax=b$ הוא $b=\begin{pmatrix}0\\
0\\
\vdots\\
0
\end{pmatrix}$, הממ״ל נקראת הומוגנית וניתן לכתוב $Ax=0$. הכוונה באגף ימין היא לוקטור האפס באורך $m$ (מספר השורות של $A$).
\end{envbox}

\begin{envbox}[thmcolor]{טענה 3.39}
א. לכל $A$ מסדר $m\times n$ ולכל $b\in\mathbb F^n$ מתקיים $\mathop{\mathrm{rank}}(A)\leq\mathop{\mathrm{rank}}(A|b)$.

ב. מתקיים $0\leq\mathop{\mathrm{rank}}(A)\leq\min\{m,n\}$ כאשר $\min\{m,n\}$ הוא המספר הקטן מבין $m,n$.
\end{envbox}

\begin{envbox}[thmcolor]{טענה 3.40}
לכל מטריצה $A$ מסדר $m\times n$ ולכל $b\in\mathbb R^n$, מספר הפתרונות לממ״ל $Ax=b$ נקבע באופן הבא:

א. אם $\mathop{\mathrm{rank}}(A)<\mathop{\mathrm{rank}}(A|b)$, אז אין לממ״ל פתרון.

ב. אם $\mathop{\mathrm{rank}}(A)=\mathop{\mathrm{rank}}(A|b)<n$, אז לממ״ל יש אינסוף פתרונות.

ג. אם $\mathop{\mathrm{rank}}(A)=\mathop{\mathrm{rank}}(A|b)=n$, אז לממ״ל יש פתרון יחיד.
\end{envbox}

\begin{envbox}[thmcolor]{טענה 3.43}
א. לכל ממ״ל הומוגנית $Ax=0$ קיים פתרון טריוויאלי שהוא וקטור האפס $x=0$, כלומר הפתרון בו כל המשתנים מתאפסים.

ב. אם $\mathop{\mathrm{rank}}(A)=n$, אז הפתרון יחיד.

ג. אם $\mathop{\mathrm{rank}}(A)<n$, אז יש אינסוף פתרונות. זה בהכרח המקרה אם מתקיים $m<n$.
\end{envbox}

\newpage

\section{פרק 4: מטריצות}

\begin{envbox}[defcolor]{הגדרה 4.2}
נגדיר את קבוצת כל המטריצות מסדר $m\times n$ מעל $\mathbb{F}$ (כאשר $\mathbb{F}=\mathbb{R}$ או $\mathbb{F}=\mathbb{C}$ ) ע״י

\[
.\mathbb{M}_{m\times n}(\mathbb{F})=\Set{\begin{pmatrix}a_{11} & a_{12} & \cdots & a_{1n}\\
a_{21} & a_{22} & \cdots & a_{2n}\\
\vdots & \vdots & \ddots & \vdots\\
a_{m1} & a_{m2} & \cdots & a_{mn}
\end{pmatrix}|\forall 1\leq i\leq m \quad \forall 1\leq j\leq n \quad a_{ij}\in \mathbb{F}}.
\]
\end{envbox}

\begin{envbox}[defcolor]{הגדרה 4.6}
עבור וקטורים $v=(a_1,...,a_n)\in\mathbb{F}^n,w=(b_1,...,b_n)\in\mathbb{F}^n$ וסקלר $c\in\mathbb{F}$ נגדיר את הפעולות הבאות:

1.  פעולת החיבור: 
\[
v+w=(a_1+b_1,...,a_n+b_n)
\]

2.  פעולת כפל בסקלר: 
\[
c \cdot v=(c\cdot a_1,\dots,c\cdot a_n)
\]
\end{envbox}

\begin{envbox}[defcolor]{הגדרה 4.7}
עבור $A,B\in\mathbb{M}_{m\times n}(\mathbb{F})$ וסקלר $c\in\mathbb{F}$ נגדיר את הפעולות הבאות:

1.  פעולת החיבור: 
\[
\begin{aligned}
    A+B&=\begin{pmatrix}a_{11} & \cdots & a_{1n}\\
    a_{21} & \cdots & a_{2n}\\
    \vdots  & \ddots & \vdots\\
    a_{m1}  & \cdots & a_{mn}
    \end{pmatrix}+\begin{pmatrix}b_{11} & \cdots & b_{1n}\\
    b_{21} & \cdots & b_{2n}\\
    \vdots & \ddots & \vdots\\
    b_{m1} & \cdots & b_{mn}
    \end{pmatrix}\\
    &=\begin{pmatrix}a_{11}+b_{11} & \cdots & a_{1n}+b_{1n} \\ a_{21}+b_{21} & \cdots & a_{2n}+b_{2n}\\
    \vdots & \ddots & \vdots\\
    a_{m1}+b_{m1} & \cdots & a_{mn}+b_{mn}
    \end{pmatrix}
    \end{aligned}
\]

2.  פעולת כפל בסקלר: 
\[
c A=\begin{pmatrix}ca_{11} & ca_{12} & \cdots & ca_{1n}\\
    ca_{21} & ca_{22} & \cdots & ca_{2n}\\
    \vdots & \vdots & \ddots & \vdots\\
    ca_{m1} & ca_{m2} & \cdots & ca_{mn}
    \end{pmatrix}
\]

    בפרט, נגדיר $-A=(-1)\cdot A$ ובהתאמה $B-A=B+(-1)\cdot A$.
\end{envbox}

\begin{envbox}[defcolor]{הגדרה 4.9}
יש דרך מקוצרת לכתיבת חיבור. במקום לכתוב את החיבור בכל רכיב, נסתכל על האיבר הכללי של כל מטריצה. נסמן את האיבר בשורה ה- $i$ ובעמודה ה- $j$ ע״י $(P)_{ij}$ לכל $P\in\mathbb{M}_{m\times n}(\mathbb{F})$. אז למעשה הגדרנו לכל $1\leq i \leq m$ ולכל $1\leq j\leq n$: 
\[
(A+B)_{ij}=(A)_{ij}+(B)_{ij}
\]

\[
(c A)_{ij}=c(A)_{ij}
\]
\end{envbox}

\begin{envbox}[defcolor]{הגדרה 4.13}
לכל $m,n\in\mathbb{N}$ נגדיר את מטריצת האפס מסדר $m\times n$ כמטריצה שכל איבריה הם $0$: 
\[
.\mathbf{0}_{m\times n}=\begin{pmatrix}0 &0 &\cdots &0 \\
0 &0 &\cdots &0 \\
\vdots & \vdots &\ddots &\vdots\\
0 &0 &\cdots &0
\end{pmatrix}
\]
\end{envbox}

\begin{envbox}[thmcolor]{טענה 4.14}
יהיו $A,B,C\in\mathbb{M}_{m\times n}(\mathbb{F})$ ו- $\alpha,\beta\in\mathbb{F}$. אז מתקיימים הכללים הבאים:

1.  חוק החילוף (קומוטטיביות): 
\[
A+B=B+A
\]

2.  חוק הקיבוץ (אסוציאטיביות): 
\[
(A+B)+C=A+(B+C)
\]

3.  חוק הקיבוץ לכפל בסקלר: 
\[
(\alpha\beta)A=\alpha(\beta A)
\]

4.  חוק הפילוג (דיסטריביוטיביות): 
\[
\alpha(A+B)=\alpha A+\alpha B
\]

5.  ניטרליות היחידה: 
\[
1\cdot A=A
\]

6.  
\[
0\cdot A=\mathbf{0}_{m\times n}
\]

7.  ניטרליות מטריצת האפס: 
\[
A+\mathbf{0}_{m\times n}=A
\]

8.  
\[
A-A=0
\]
\end{envbox}

\begin{envbox}[defcolor]{הגדרה 4.16}
בהינתן $a_1,a_2,...,a_n\in\mathbb{F}$ נסמן את סכומם בעזרת כתיבה עם $\Sigma$ ואינדקס $l$.

\[
.\sum_{l=1}^n a_l=a_1+a_2+a_3+\cdots+a_n.
\]

האינדקס $l$ רץ על פני הקבוצה $\Set{1,2,3,\cdots,n}$ וכל הצבה תורמת איבר לסכום.
\end{envbox}

\begin{envbox}[defcolor]{הגדרה 4.18}
בהינתן $A\in\mathbb{M}_{m\times n}(\mathbb{F})$ ו- 
\[
,x=\begin{pmatrix}
x_1\\
x_2\\
\vdots \\
x_n
\end{pmatrix}\in\mathbb{F}^n
\]
 נגדיר את המכפלה 
\[
Ax=\begin{pmatrix}
y_1\\
y_2\\
\vdots \\
y_m
\end{pmatrix}\in\mathbb{F}^m
\]
 עם קוארדינטות הנתונות ע״י 
\[
y_i=\sum_{j=1}^n (A)_{ij}x_j.
\]
 לכל $1\leq i\leq m$. נסמן $y_i=(Ax)_i$ עבור הקוארדינטות.
\end{envbox}

\begin{envbox}[thmcolor]{טענה 4.23}
יהיו $A,B\in\mathbb{M}_{m\times n}(\mathbb{F})$, $v,w\in\mathbb{F}^n$ ו- $\alpha\in\mathbb{F}$.

1.  
\[
A(v+w)=Av+Aw
\]

2.  
\[
(A+B)v=Av+Bv
\]

3.  
\[
A(\alpha v)=\alpha(Av)
\]
\end{envbox}

\begin{envbox}[thmcolor]{טענה 4.24}
תהי $Ax=0$ ממ״ל הומוגנית.

1.  קבוצת הפתרונות סגורה לחיבור, כלומר לכל שני פתרונות $v,w\in\mathbb{F}^n$ גם $v+w$ הוא פתרון.

2.  קבוצת הפתרונות סגורה לכפל בסקלר, כלומר לכל פתרון $v\in\mathbb{F}^n$ ולכל $\alpha\in\mathbb{F}$ גם $\alpha v$ הוא פתרון.
\end{envbox}

\begin{envbox}[thmcolor]{טענה 4.25}
תהי $Ax=b$ ממ״ל לא הומוגנית. נניח שקיים לה פתרון (לא בהכרח יחיד) $v_0\in\mathbb{F}^n$. אז מתקיים הקשר הבא בין קבוצת הפתרונות של הממ״ל הלא הומוגנית לקבוצת הפתרונות של הממ״ל ההומוגנית $Ax=0$:

\[
.\Set{x\in\mathbb{F}^n|Ax=b}=\Set{v_0+y|y\in\mathbb{F}^n,Ay=0}.
\]
\end{envbox}

\begin{envbox}[defcolor]{הגדרה 4.27}
בהינתן $A\in\mathbb{M}_{m\times k}(\mathbb{F}),B\in\mathbb{M}_{k\times n}(\mathbb{F})$ נגדיר את מטריצת הכפל $AB$ ע״י 
\[
(AB)_{ij}=\sum_{l=1}^k (A)_{il}(B)_{lj}=(A)_{i1}(B)_{1j}+(A)_{i2}(B)_{2j}+...+(A)_{ik}(B)_{kj}.
\]
 לכל $1\leq i\leq m, 1\leq j\leq n$.
\end{envbox}

\begin{envbox}[thmcolor]{טענה 4.32}
יהיו $A\in\mathbb{M}_{m\times k}(\mathbb{F}), B\in\mathbb{M}_{k\times n}(\mathbb{F})$. נניח שהעמודות של $B$ הם הוקטורים $v_1,v_2,...,v_n\in\mathbb{F}^k$. אז מתקיים

\[
AB=A\begin{pmatrix}
| & | & & | \\
v_1 & v_2 & \cdots & v_n \\
| & | & & |
\end{pmatrix}=\begin{pmatrix}
| & | & & | \\
Av_1 & Av_2 & \cdots & Av_n \\
| & | & & |
\end{pmatrix}
\]

כלומר התוצאה של כפל ב- $A$ משמאל הוא שכל וקטור עמודה מוכפל ב- $A$ לחוד.
\end{envbox}

\begin{envbox}[thmcolor]{טענה 4.33}
יהיו $A_1,A_2\in\mathbb{M}_{m\times n}(\mathbb{F})$, $B_1,B_2\in\mathbb{M}_{n\times p}(\mathbb{F})$, $C\in\mathbb{M}_{p\times q}(\mathbb{F})$ ו- $\alpha\in\mathbb{F}$. אז מתקיים

1.  
\[
(A_1+A_2)B_1=A_1B_1+A_2B_1
\]

2.  
\[
A_1(B_1+B_2)=A_1B_1+A_1B_2
\]

3.  
\[
A_1(B_1C)=(A_1B_1)C
\]

4.  
\[
(\alpha A_1)B_1=\alpha(A_1B_1)=A_1(\alpha B_1)
\]

5.  
\[
\mathbf{0}_{m\times n} B_1=\mathbf{0}_{m\times p}=A_1\mathbf{0}_{n\times p}
\]
\end{envbox}

\begin{envbox}[defcolor]{הגדרה 4.34}
מטריצת היחידה מסדר $n\times n$ (אפשר גם לומר מסדר $n$) מוגדרת להיות 
\[
.I_n=I_{n\times n}=.
\begin{pmatrix}
1 &0 &\cdots &0\\
0 &1 &\cdots &0\\
\vdots &\vdots &\ddots &\vdots\\
0 &0 &\cdots &1
\end{pmatrix}
\]

זוהי המטריצה שבה כל איבר על האלכסון הראשי הוא $1$, וכל שאר האיברים הם $0$. או באופן מתמטי:

\[
\begin{aligned}
\forall 1\leq i\leq n \quad (I_n)_{ii}=1 \\
\forall 1\leq i,j\leq n  \quad i\neq j \implies &(I_n)_{ij}=0
\end{aligned}
\]
\end{envbox}

\begin{envbox}[thmcolor]{טענה 4.36}
לכל $A\in\mathbb{M}_{m\times n}(\mathbb{F})$ מתקיים $AI_n=A=I_m A$.
\end{envbox}

\begin{envbox}[defcolor]{הגדרה 4.38}
תהי $A\in\mathbb{M}_{m\times n}(\mathbb{F})$. נגדיר את המטריצה המשוחלפת $A^t\in\mathbb{M}_{n\times m}(\mathbb{F})$ ע״י 
\[
(A^t)_{ij}=A_{ji}
\]
 לכל $1\leq i\leq n,1\leq j\leq m$.
\end{envbox}

\begin{envbox}[thmcolor]{טענה 4.41}
לכל $A,B\in\mathbb{M}_{m\times n}(\mathbb{F}), C\in\mathbb{M}_{n\times k}(\mathbb{F})$ ולכל $\alpha\in\mathbb{F}$ מתקיים:

1.  
\[
(A^t)^t=A
\]

2.  
\[
(A+B)^t=A^t+B^t
\]

3.  
\[
(\alpha A)^t=\alpha A^t
\]

4.  
\[
(AC)^t=C^tA^t
\]
\end{envbox}

\begin{envbox}[defcolor]{הגדרה 4.45}
מטריצה $A$ נקראת ריבועית אם מספר השורות שלה שווה למספר העמודות, כלומר אם קיים $n\in\mathbb{N}$ עבורו $A\in\mathbb{M}_{n\times n}(\mathbb{F})$.
\end{envbox}

\begin{envbox}[defcolor]{הגדרה 4.48}
תהי $A\in\mathbb{M}_{n\times n}(\mathbb{F})$. אז נגדיר

\[
A^0=I_n, \quad A^1=A, \quad A^2=A\cdot A, \quad A^3=A\cdot A\cdot A
\]

ולכל $k\in\mathbb{N}$

\[
.A^k = \underbrace{A \cdot A \cdot \ldots \cdot A}_{\text{פעמים} \ k}
\]

עבור פולינום 
\[
p(x)=a_mx^m+a_{m-1}x^{m-1}+...+a_1x+a_0
\]
 נגדיר את ההצבה של $A$ בו ע״י 
\[
.p(A)=a_mA^m+a_{m-1}A^{m-1}+...+a_1A+a_0I_n.
\]
\end{envbox}

\begin{envbox}[defcolor]{הגדרה 4.51}
1.  $A$ נקראת סימטרית אם מתקיים $A^t=A$, או באופן שקול: 
\[
\forall 1\leq i,j\leq n \quad (A)_{ji}=(A)_{ij}
\]

2.  נקראת אנטי-סימטרית אם מתקיים $A^t=-A$, או באופן שקול: 
\[
\forall 1\leq i,j\leq n \quad (A)_{ji}=-(A)_{ij}
\]
\end{envbox}

\begin{envbox}[thmcolor]{טענה 4.53}
תהי $P\in\mathbb{M}_{n\times n}(\mathbb{F})$ מטריצה ריבועית.

1.  אם $P$ אנטי-סימטרית, אז כל איבריה על האלכסון הראשי שווים ל- $0$.

2.  אם $P$ סימטרית וגם אנטי-סימטרית, אז $P=0$.

3.  לכל $B\in\mathbb{M}_{n\times n}(\mathbb{F})$ קיימות $S\in\mathbb{M}_{n\times n}(\mathbb{F})$ סימטרית ו- $A\in\mathbb{M}_{n\times n}(\mathbb{F})$ אנטי-סימטרית כך שמתקיים $B=S+A$.
\end{envbox}

\begin{envbox}[defcolor]{הגדרה 4.57}
מטריצה $A\in\mathbb{M}_{n\times n}(\mathbb{F})$ נקראת:

1.  אלכסונית אם לכל $i\neq j$ מתקיים $(A)_{ij}=0$, כלומר כל האיברים מחוץ לאלכסון הראשי שווים ל- $0$.

2.  משולשית עליונה אם לכל $i>j$ מתקיים $(A)_{ij}=0$, כלומר כל האיברים מתחת לאלכסון הראשי שווים ל- $0$.

3.  משולשית תחתונה אם לכל $i<j$ מתקיים $(A)_{ij}=0$, כלומר כל האיברים מעל לאלכסון הראשי שווים ל- $0$.
\end{envbox}

\begin{envbox}[thmcolor]{טענה 4.60}
יהיו $A,B\in\mathbb{M}_{n\times n}(\mathbb{F})$.

1.  אם $A,B$ משולשיות עליונות, אז $AB$ גם משולשית עליונה.

2.  אם $A,B$ משולשיות תחתונות, אז $AB$ גם משולשית תחתונה.

3.  אם $A,B$ אלכסוניות, אז $AB$ גם אלכסונית.
\end{envbox}

\begin{envbox}[defcolor]{הגדרה 4.62}
מטריצה $A\in\mathbb{M}_{n\times n}(\mathbb{F})$ נקראת סקלרית אם קיים $\alpha\in\mathbb{F}$ כך ש- $A=\alpha I_n$.
\end{envbox}

\begin{envbox}[defcolor]{הגדרה 4.64}
מטריצה $A\in\mathbb{M}_{n\times n}(\mathbb{F})$ נקראת אלמנטרית אם היא מתקבלת ממטריצת היחידה $I_n$ ע״י פש״א אחת.

עבור פש״א $S$, נסמן ב- $S(I_n)$ את המטריצה האלמנטרית המתאימה שמתקבלת מ- $I_n$ ע״י $S$. באופן כללי, לכל מטריצה $A$ נסמן ב- $S(A)$ את המטריצה שמתקבלת מ- $A$ ע״י $S$.
\end{envbox}

\begin{envbox}[thmcolor]{טענה 4.66}
יהיו $A\in\mathbb{M}_{m\times n}(\mathbb{F})$ ו- $S$ פש״א. אז מתקיים $S(I_m)A=S(A)$.
\end{envbox}

\begin{envbox}[thmcolor]{מסקנה 4.68}
יהיו $A,A'\in\mathbb{M}_{m\times n}(\mathbb{F})$ שקולות שורה. אז קיימת $B\in\mathbb{M}_{m\times m}(\mathbb{F})$ כך ש- $A'=BA$.
\end{envbox}

\begin{envbox}[defcolor]{הגדרה 4.70}
נגדיר את וקטורי היחידה $e_1,e_2,\cdots,e_n\in\mathbb{F}^n$ ע״י 
\[
e_i=\begin{pmatrix}
0 \\
0 \\
\vdots \\
0 \\
1 \\
0\\
\vdots \\
0
\end{pmatrix}
\]
 כאשר כל הקוארדינטות הן $0$ חוץ מהקוארדינטה ה- $i$ השווה ל- $1$. נציין שאלה וקטורי העמודה של $I_n$, וגם וקטורי השורה עד כדי שחלוף (ההבדל בין שורה לעמודה).
\end{envbox}

\begin{envbox}[thmcolor]{טענה 4.72}
אם $A'\in\mathbb{M}_{n\times n}(\mathbb{F})$ מדורגת קנונית, אז התנאים הבאים שקולים:

1.  $A'=I_n$

2.  $\mathop{\mathrm{rank}}(A')=n$

3.  לממ״ל $A'x=0$ יש פתרון יחיד.

4.  לכל $b'\in\mathbb{F}^n$ יש פתרון יחיד לממ״ל $A'x=b'$.

5.  לכל $b'\in\mathbb{F}^n$ יש פתרון לממ״ל $A'x=b'$.
\end{envbox}

\begin{envbox}[thmcolor]{טענה 4.73}
לכל $A\in\mathbb{M}_{n\times n}(\mathbb{F})$ התנאים הבאים שקולים:

1.  $A$ שקולת שורה ל- $I_n$.

2.  $\mathop{\mathrm{rank}}(A)=n$

3.  לממ״ל $Ax=0$ יש פתרון יחיד.

4.  לכל $b\in\mathbb{F}^n$ יש פתרון יחיד לממ״ל $Ax=b$.

5.  לכל $b\in\mathbb{F}^n$ יש פתרון לממ״ל $Ax=b$.
\end{envbox}

\begin{envbox}[defcolor]{הגדרה 4.74}
מטריצה $A\in\mathbb{M}_{n\times n}(\mathbb{F})$ נקראת הפיכה משמאל אם קיימת $B\in\mathbb{M}_{n\times n}(\mathbb{F})$ כך ש- $BA=I_n$. $A$ נקראת הפיכה מימין אם קיימת $C\in\mathbb{M}_{n\times n}(\mathbb{F})$ כך ש- $AC=I_n$. אם שני התנאים מתקיימים, $A$ נקראת הפיכה.
\end{envbox}

\begin{envbox}[thmcolor]{משפט 4.75}
לכל $A\in\mathbb{M}_{n\times n}(\mathbb{F})$ התנאים הבאים שקולים:

1.  $A$ שקולת שורה ל- $I_n$.

2.  $\mathop{\mathrm{rank}}(A)=n$

3.  לממ״ל $Ax=0$ יש פתרון יחיד.

4.  לכל $b\in\mathbb{F}^n$ יש פתרון יחיד לממ״ל $Ax=b$.

5.  לכל $b\in\mathbb{F}^n$ יש פתרון לממ״ל $Ax=b$.

6.  $A$ הפיכה משמאל.

7.  $A$ הפיכה מימין.

8.  $A$ הפיכה.
\end{envbox}

\begin{envbox}[thmcolor]{טענה 4.76}
יהיו $A,B,C\in\mathbb{M}_{n\times n}(\mathbb{F})$ כך ש- $BA=I_n=AC$. אז $B=C$.
\end{envbox}

\begin{envbox}[thmcolor]{מסקנה 4.77}
יהיו $A,B\in\mathbb{M}_{n\times n}(\mathbb{F})$. אז מתקיים 
\[
.AB=I_n \iff BA=I_n
\]
\end{envbox}

\begin{envbox}[thmcolor]{מסקנה 4.78}
יהיו $A,B_1,B_2\in\mathbb{M}_{n\times n}(\mathbb{F})$ כך ש- $B_1A=B_2A=I_n$. אז $B_1=B_2$.
\end{envbox}

\begin{envbox}[defcolor]{הגדרה 4.79}
תהי $A\in\mathbb{M}_{n\times n}(\mathbb{F})$ הפיכה. נסמן ב- $A^{-1}$ את המטריצה שמקיימת 
\[
.AA^{-1}=A^{-1}A=I_n
\]
 נקרא לה המטריצה ההופכית ל- $A$.
\end{envbox}

\begin{envbox}[thmcolor]{מסקנה 4.84}
1.  תהי $A\in\mathbb{M}_{n\times n}(\mathbb{F})$ הפיכה. אז מתקיים 
\[
.A^{-1}=E_kE_{k-1}\cdots E_1.
\]
 כאשר $E_1,E_2,...,E_k$ הן המטריצות האלמנטריות המתאימות לפש״אות $S_1,S_2,...,S_k$ בדירוג הקנוני של $A$.

2.  $A$ היא מכפלה של מטריצות אלמנטריות: 
\[
.A=E_1^{-1}E_2^{-1}\cdots E_k^{-1}.
\]
\end{envbox}

\newpage

\section{פרק 5: דטרמיננטה}

\begin{envbox}[defcolor]{הגדרה 5.1}
עבור $A\in\mathbb{M}_{n\times n}(\mathbb{F})$ נגדיר את $\det(A)$ ע״י התכונות הבאות:

1.  $.\det(I_n)=1$

2.  אם $A\xrightarrow{\alpha R_i\to R_i}A'$ עבור $1\leq i \leq n$, אז $\det(A')=\alpha \det(A)$.

3.  אם $A\xrightarrow{R_i+\alpha R_j\to R_i}A'$ עבור $i\neq j$, אז $\det(A')=\det(A)$.

4.  אם $A\xrightarrow{R_i\leftrightarrow R_j}A'$ עבור $i\neq j$, אז $\det(A')=-\det(A)$.
\end{envbox}

\begin{envbox}[thmcolor]{טענה 5.5}
יהיו $A\in\mathbb{M}_{n\times n}(\mathbb{F})$ ו- $\alpha\in\mathbb{F}$. אז מתקיים 
\[
.\det(\alpha A)=\alpha^n\det(A)
\]
\end{envbox}

\begin{envbox}[thmcolor]{טענה 5.6}
אם למטריצה $A\in\mathbb{M}_{n\times n}(\mathbb{F})$ יש שורת אפסים, אז $\det(A)=0$.
\end{envbox}

\begin{envbox}[thmcolor]{טענה 5.7}
יהיו $D,U\in\mathbb{M}_{n\times n}(\mathbb{F})$.

1.  אם $D$ אלכסונית מהצורה

    
\[
,D=\begin{pmatrix}
    d_1 &0 & \cdots & 0 \\
    0 & d_2 & \cdots & 0 \\
    \vdots & \vdots & \ddots & 0 \\
    0 & 0 & \cdots & d_n
    \end{pmatrix}
\]

    אז מתקיים

    
\[
.\det(D)=d_1d_2\cdots d_n
\]

2.  אם $U$ משולשית עליונה מהצורה

    
\[
U=\begin{pmatrix}
    d_1 &* & \cdots & * \\
    0 & d_2 & \cdots & * \\
    \vdots & \vdots & \ddots & * \\
    0 & 0 & \cdots & d_n
    \end{pmatrix}
\]

    כאשר הכוכביות מציינות איברים כלשהם ובנוסף $d_i\neq 0$ לכל $1\leq i\leq n$, אז מתקיים

    
\[
.\det(U)=d_1d_2\cdots d_n
\]
\end{envbox}

\begin{envbox}[thmcolor]{טענה 5.10}
`\label{prop:det-order-2}`{=latex} עבור $n=2$ הפונקציה $\det:\mathbb{M}_{2\times 2}(\mathbb{F})\to\mathbb{F}$ מהגדרה [7.1](\#def-det) מקיימת:

1.  לכל 
\[
A=\begin{pmatrix}
    a & b \\
    c & d
    \end{pmatrix}
\]
 מתקיים 
\[
.\det(A)=ad-bc
\]

2.  $A$ הפיכה אם ורק אם $\det(A)\neq 0$. במקרה זה מתקיים 
\[
.A^{-1}=\frac{1}{ad-bc}\begin{pmatrix}
    d & -b \\
    -c & a
    \end{pmatrix}
\]
\end{envbox}

\begin{envbox}[thmcolor]{טענה 5.14}
`\label{prop:elementary-det}`{=latex} יהיו $E,A\in\mathbb{M}_{n\times n}(\mathbb{F})$ כך ש- $E$ מטריצה אלמנטרית המתאימה לפש״א $S$, כלומר $E=S(I_n)$. אז מתקיים 
\[
.\det(EA)=\det(E)\det(A)
\]
\end{envbox}

\begin{envbox}[thmcolor]{משפט 5.15}
תהי $A\in\mathbb{M}_{n\times n}(F)$. אז $A$ הפיכה אם ורק אם $\det(A)\neq 0$.
\end{envbox}

\begin{envbox}[thmcolor]{טענה 5.17}
`\label{prop:det-product}`{=latex} יהיו $A,B\in\mathbb{M}_{n\times n}(\mathbb{F})$. אז מתקיים 
\[
.\det(AB)=\det(A)\det(B)
\]
\end{envbox}

\begin{envbox}[thmcolor]{מסקנה 5.18}
`\label{cor:det-inverse}`{=latex} תהי $A\in\mathbb{M}_{n\times n}(\mathbb{F})$ הפיכה. אז מתקיים 
\[
.\det(A^{-1})=\frac{1}{\det(A)}
\]
\end{envbox}

\begin{envbox}[thmcolor]{טענה 5.21}
`\label{prop:det-transpose}`{=latex} תהי $A\in\mathbb{M}_{n\times n}(\mathbb{F})$. אז מתקיים 
\[
.\det(A^t)=\det(A)
\]
\end{envbox}

\begin{envbox}[thmcolor]{מסקנה 5.22}
תהי $U,L\in\mathbb{M}_{n\times n}(\mathbb{F})$.

1.  אם $U$ משולשית עליונה מהצורה

    
\[
U=\begin{pmatrix}
    d_1 &* & \cdots & * \\
    0 & d_2 & \cdots & * \\
    \vdots & \vdots & \ddots & * \\
    0 & 0 & \cdots & d_n
    \end{pmatrix}
\]

    כאשר הכוכביות מציינות איברים כלשהם, אז מתקיים

    
\[
.\det(U)=d_1d_2\cdots d_n
\]

2.  אם $L$ משולשית תחתונה מהצורה

    
\[
L=\begin{pmatrix}
    d_1 &0 & \cdots & 0 \\
    * & d_2 & \cdots & 0 \\
    \vdots & \vdots & \ddots & 0 \\
    * & * & \cdots & d_n
    \end{pmatrix}
\]

    כאשר הכוכביות מציינות איברים כלשהם, אז מתקיים

    
\[
.\det(L)=d_1d_2\cdots d_n
\]
\end{envbox}

\begin{envbox}[defcolor]{הגדרה 5.24}
תהי $A\in\mathbb{M}_{n\times n}(\mathbb{F})$. נגדיר את המינור $M_{i,j}(A)$ להיות המטריצה שמתקבלת מ- $A$ ע״י מחיקת השורה ה- $i$ והעמודה ה- $j$.
\end{envbox}

\begin{envbox}[thmcolor]{טענה 5.27}
`\label{prop:laplace}`{=latex} תהי 
\[
.A=\begin{pmatrix}a_{11} & \cdots & a_{1n}\\
\vdots & \ddots & \vdots\\
a_{n1} & \cdots & a_{nn}
\end{pmatrix}
\]
 אז מתקיים:

1.  לכל $1\leq i \leq n$, ניתן לפתח את הדטרמיננטה לפי השורה ה- $i$: 
\[
\det(A)=\sum_{j=1}^n (-1)^{i+j}a_{ij}\det(M_{i,j}(A))
\]

2.  לכל $1\leq j \leq n$, ניתן לפתח את הדטרמיננטה לפי העמודה ה- $j$: 
\[
\det(A)=\sum_{i=1}^n (-1)^{i+j}a_{ij}\det(M_{i,j}(A))
\]
\end{envbox}

\newpage

\section{פרק 6: מרחבים וקטוריים}

\begin{envbox}[defcolor]{הגדרה 6.9}
יהי $V$ מ״ו מעל $\mathbb{F}$. תת-קבוצה $W\subseteq V$ תיקרא תת-מרחב של $V$ אם $W$ היא מ״ו מעל $\mathbb{F}$ בפני עצמה, עם פעולות החיבור והכפל בסקלר מ- $V$.
\end{envbox}

\begin{envbox}[defcolor]{הגדרה 6.14}
תהי $A\in\mathbb{M}_{m\times n}(\mathbb{F})$. קבוצת הפתרונות לממ״ל ההומוגנית $Ax=\underline{0}$ נקראת מרחב הפתרונות של $A$ ומסומנת ע״י $\mathop{\mathrm{N}}(A)$.
\end{envbox}

\begin{envbox}[thmcolor]{טענה 6.15}
$\mathop{\mathrm{N}}(A)$ הוא תת-מרחב של $\mathbb{F}^n$.
\end{envbox}

\begin{envbox}[defcolor]{הגדרה 6.25}
יהי $V$ מ״ו מעל $\mathbb{F}$, ויהיו $v_1,...,v_k\in V$. צירוף לינארי (צ״ל) של $v_1,...,v_k$ הוא וקטור מהצורה 
\[
v=\alpha_1v_1+...+\alpha_kv_k
\]
 כאשר $\alpha_1,...,\alpha_k\in\mathbb{F}$.
\end{envbox}

\begin{envbox}[defcolor]{הגדרה 6.28}
יהי $V$ מ״ו מעל $\mathbb{F}$, ויהיו $v_1,v_2,...,v_k\in V$. אז נגדיר את קבוצת כל הצירופים הלינאריים של וקטורים אלה: 
\[
\mathop{\mathrm{Span}}(v_1,...,v_k)=\Set{\alpha_1v_1+...+\alpha_kv_k|\alpha_1,...,\alpha_k\in\mathbb{F}}
\]
\end{envbox}

\begin{envbox}[defcolor]{הגדרה 6.31}
יהי $V$ מ״ו מעל $\mathbb{F}$. נאמר ש- $V$ נוצר סופית אם קיימים $v_1,...,v_k\in V$ כך ש- $V=\mathop{\mathrm{Span}}(v_1,...,v_k)$. במקרה זה גם נאמר שהקבוצה $\Set{v_1,...,v_k}$ פורשת את $V$. ניתן גם לומר (באופן פחות פורמלי) שהוקטורים $v_1,...,v_k$ פורשים את $V$.
\end{envbox}

\begin{envbox}[defcolor]{הגדרה 6.36}
יהי $V$ מ״ו מעל $\mathbb{F}$. קבוצת וקטורים $\Set{v_1,...,v_k}\subseteq V$ נקראת:

-   תלויה לינארית (ת״ל) אם קיימים $\alpha_1,...,\alpha_k\in\mathbb{F}$ שאינם כולם שווים ל- $0$ כך ש- 
\[
.\alpha_1v_1+...+\alpha_kv_k=0_V
\]

    במקרה זה גם נאמר שהצירוף הלינארי באגף שמאל הוא לא טריוויאלי.

-   בלתי תלויה לינארית (בת״ל) אם לכל $\alpha_1,...,\alpha_k\in\mathbb{F}$ עבורם מתקיים 
\[
,\alpha_1v_1+...+\alpha_kv_k=0_V
\]
 בהכרח מתקיים $\alpha_1=...=\alpha_k=0$. צירוף לינארי זה נקרא טריוויאלי.
\end{envbox}

\begin{envbox}[thmcolor]{מסקנה 6.42}
יהי $V$ מ״ו מעל $\mathbb{F}$, ויהיו $v_1,...,v_k,v_{k+1},...,v_l\in V$.

א. אם $\Set{v_1,...,v_k}$ ת״ל, אז $\Set{v_1,...,v_k,v_{k+1},...,v_l}$ ת״ל.

ב. אם $\Set{v_1,...,v_k,v_{k+1},...,v_l}$ בת״ל, אז $\Set{v_1,...,v_k}$ בת״ל.
\end{envbox}

\begin{envbox}[thmcolor]{מסקנה 6.44}
יהי $V$ מ״ו מעל $\mathbb{F}$, ויהיו $\Set{v_1,...,v_k}\in V$ בת״ל ו- $w\in V$ כך שמתקיים $w\notin\mathop{\mathrm{Span}}(v_1,...,v_k)$. אז $\Set{v_1,...,v_k,w}$ בת״ל.
\end{envbox}

\begin{envbox}[thmcolor]{מסקנה 6.45}
`\label{cor:basis-extraction}`{=latex} אם $\Set{v_1,...,v_k}$ קבוצה הפורשת מ״ו $V$ מעל $\mathbb{F}$, אז תת-הקבוצה 
\[
\Set{v_i|1\leq i\leq k, v_i\notin\mathop{\mathrm{Span}}(v_1,...,v_{i-1})}
\]
 היא בת״ל ופורשת את $V$.
\end{envbox}

\begin{envbox}[defcolor]{הגדרה 6.48}
יהי $V$ מ״ו מעל $\mathbb{F}$. תת-קבוצה $B=\Set{v_1,...,v_n}\subseteq V$ נקראת בסיס ל- $V$ אם היא בת״ל וגם פורשת את $V$, כלומר $v_1,...,v_n$ בת״ל ובנוסף $\mathop{\mathrm{Span}}(B)=V$.
\end{envbox}

\begin{envbox}[defcolor]{הגדרה 6.51}
יהי $V=\Set{0_V}$. עבור קבוצת וקטורים ריקה $\emptyset=\Set{ }$ נגדיר 
\[
.\mathop{\mathrm{Span}}(\emptyset)=\Set{0_V}
\]
\end{envbox}

\begin{envbox}[thmcolor]{טענה 6.53}
יהי $V$ מ״ו מעל $\mathbb{F}$. אז התנאים הבאים שקולים:

-   ל- $V$ קיים בסיס.

-   $V$ נוצר סופית, כלומר קיימת תת-קבוצה סופית הפורשת אותו.
\end{envbox}

\begin{envbox}[thmcolor]{מסקנה 6.55}
`\label{cor:size-dim}`{=latex} יהיו $v_1,...,v_k\in\mathbb{F}^n$.

א. אם $\Set{v_1,...,v_k}$ בת״ל, אז $k\leq n$.

ב. אם $\Set{v_1,...,v_k}$ פורשת את $\mathbb{F}^n$, אז $k\geq n$.

ג. אם $\Set{v_1,...,v_k}$ היא בסיס ל- $\mathbb{F}^n$, אז $k=n$.
\end{envbox}

\begin{envbox}[thmcolor]{מסקנה 6.59}
`\label{cor:row-equivalence}`{=latex} יהיו $A,B\in\mathbb{M}_{m\times n}(\mathbb{F})$ שקולות שורה. אז $\mathop{\mathrm{Row}}(A)=\mathop{\mathrm{Row}}(B)$.
\end{envbox}

\begin{envbox}[defcolor]{הגדרה 6.62}
בסיס סדור $B=\Set{v_1,...,v_n}$ למ״ו $V$ הוא בסיס יחד עם סדר על איבריו (כלומר יש חשיבות לסדר כתיבתם).
\end{envbox}

\begin{envbox}[defcolor]{הגדרה 6.65}
יהי $V$ מ״ו מעל $\mathbb{F}$, ויהי $B=\Set{v_1,...,v_n}$ בסיס סדור ל- $V$. לכל $v\in V$ קיימים $\alpha_1,...,\alpha_n\in\mathbb{F}$ יחידים כך ש- $v=\alpha_1v_1+...+\alpha_nv_n$. אז עבור $v$ נגדיר את וקטור הקוארדינטות שלו ביחס ל- $B$ ע״י 
\[
.[v]_B=\begin{pmatrix}
\alpha_1 \\
\vdots \\
\alpha_n
\end{pmatrix}\in\mathbb{F}^n
\]
\end{envbox}

\begin{envbox}[thmcolor]{מסקנה 6.68}
`\label{cor:span-coordinates}`{=latex} יהיו $V$ מ״ו מעל $\mathbb{F}$, $B=\Set{v_1,...,v_n}$ בסיס סדור ל- $V$, ו- $u,w_1,...,w_k\in V$. אז מתקיים 
\[
.u\in\mathop{\mathrm{Span}}(w_1,...,w_k) \iff [u]_B\in\mathop{\mathrm{Span}}([w_1]_B,...,[w_k]_B)
\]
\end{envbox}

\begin{envbox}[thmcolor]{מסקנה 6.71}
`\label{cor:size-dim2}`{=latex} יהיו $V$ מ״ו מעל $\mathbb{F}$, $B=\Set{b_1,...,b_n}$ בסיס ל- $V$ ו- $v_1,...,v_k\in V$.

א. אם $\Set{v_1,...,v_k}$ בת״ל, אז $k\leq n$.

ב. אם $\Set{v_1,...,v_k}$ פורשת את $V$, אז $k\geq n$.

ג. אם $\Set{v_1,...,v_k}$ היא בסיס ל- $V$, אז $k=n$.
\end{envbox}

\begin{envbox}[thmcolor]{מסקנה 6.72}
יהי $V$ מ״ו נוצר סופית, כלומר קיים לו בסיס. אז כל שני בסיסים ל- $V$ הם שווי גודל.
\end{envbox}

\begin{envbox}[defcolor]{הגדרה 6.73}
גודל בסיס כלשהו של מ״ו $V$ נקרא המימד של $V$. סימונו הוא $\dim(V)$.
\end{envbox}

\begin{envbox}[thmcolor]{מסקנה 6.78}
`\label{cor:subspace-dim}`{=latex} יהי $V$ מ״ו נוצר סופית מעל $\mathbb{F}$, ויהי $W$ תת-מרחב. אז מתקיים:

א. $W$ נוצר סופית ומתקיים $\dim(W)\leq \dim(V)$.

ב. אם $\dim(W)=\dim(V)$, אז $W=V$.
\end{envbox}

\newpage

\section{פרק 7: תת-מרחבים}

\begin{envbox}[defcolor]{הגדרה 7.1}
יהי $V$ מ״ו מעל $\mathbb{F}$, ויהיו $W_1,W_2\subseteq V$ תת-מרחבים. אז החיתוך $W_1\cap W_2$ הוא קבוצת כל הוקטורים השייכים גם ל- $W_1$ וגם ל- $W_2$, כלומר 
\[
.W_1\cap W_2=\Set{v\in V|v\in W_1,v\in W_2}
\]
\end{envbox}

\begin{envbox}[thmcolor]{טענה 7.4}
יהי $V$ מ״ו מעל $\mathbb{F}$, ויהיו $W_1,W_2\subseteq V$ תת-מרחבים. אז $W_1\cap W_2$ הוא תת-מרחב של $V$.
\end{envbox}

\begin{envbox}[defcolor]{הגדרה 7.8}
יהי $V$ מ״ו מעל $\mathbb{F}$, ויהיו $W_1,W_2\subseteq V$ תת-מרחבים. אז הסכום $W_1+W_2$ מוגדר ע״י 
\[
.W_1+W_2=\Set{w_1+w_2|w_1\in W_1,w_2\in W_2}
\]
\end{envbox}

\begin{envbox}[thmcolor]{טענה 7.9}
יהי $V$ מ״ו מעל $\mathbb{F}$, ויהיו $W_1,W_2\subseteq V$ תת-מרחבים. אז $W_1+W_2$ הוא תת-מרחב של $V$, המקיים $W_1\subseteq W_1+W_2$ וגם $W_2\subseteq W_1+W_2$.
\end{envbox}

\begin{envbox}[thmcolor]{מסקנה 7.14}
יהיו $V$ מ״ו מעל $\mathbb{F}$ ו- $W_1,W_2\subseteq V$. אז מתקיים: 
\[
\begin{aligned}
\dim(W_1\cap W_2)\leq \dim(W_1)\leq\dim(W_1+W_2)\leq\dim(V)
\end{aligned}
\]
 ובאופן דומה: 
\[
\dim(W_1\cap W_2)\leq \dim(W_2)\leq\dim(W_1+W_2)\leq\dim(V)
\]
\end{envbox}

\begin{envbox}[thmcolor]{משפט 7.16}
`\label{thm:dim-formula}`{=latex} יהי $V$ מ״ו נוצר סופית מעל $\mathbb{F}$, ויהיו $W_1,W_2\subseteq V$ תת-מרחבים. אז מתקיים 
\[
.\dim(W_1+W_2)=\dim(W_1)+\dim(W_2)-\dim(W_1\cap W_2)
\]
\end{envbox}

\newpage


\end{document}
